\documentclass[11pt]{article}
\usepackage{amsmath,amssymb,amsthm}
\usepackage[margin=1in]{geometry}
\usepackage{array}
\usepackage{tabularx}
\usepackage{booktabs}
\usepackage[hidelinks]{hyperref}

\theoremstyle{plain}
\newtheorem{theorem}{Theorem}[section]
\newtheorem{lemma}[theorem]{Lemma}
\newtheorem{proposition}[theorem]{Proposition}
\newtheorem{corollary}[theorem]{Corollary}
\theoremstyle{definition}
\newtheorem{definition}[theorem]{Definition}
\newtheorem{example}[theorem]{Example}
\newtheorem{fact}[theorem]{Fact}
\theoremstyle{remark}
\newtheorem{remark}[theorem]{Remark}

\newcommand{\Z}{\mathbb{Z}}
\newcommand{\sm}{\setminus}
\DeclareMathOperator{\exc}{exc}
\DeclareMathOperator{\solo}{solo}
\newcommand{\Adm}{\mathrm{Adm}}
\newcommand{\dg}{\operatorname{deg}}

\newcommand{\partheading}[1]{%
  \par\vspace{1.6\baselineskip}%
  \noindent\rule{\linewidth}{0.8pt}\par\vspace{0.35\baselineskip}%
  \begingroup\centering\normalsize\scshape #1\par\endgroup
  \vspace{0.35\baselineskip}\noindent\rule{\linewidth}{0.8pt}%
  \par\vspace{1.0\baselineskip}}

\title{Extreme endpoints of low-excess polyhexes are never articulation cells}
\author{H. Utsunomiya}
\date{}
\begin{document}
\maketitle
\begin{abstract}
Let \(S\) be a polyhex with at least four cells whose edge perimeter exceeds the minimum
for its size by at most two. We prove that no endpoint of any of the six extreme lines of
\(S\) is an articulation cell of \(S\). The proof combines a local structure argument at
an extreme-line endpoint with a bridge-excess lemma stating that a degree-two articulation
cell forces excess at least four. Both hypotheses are shown to be necessary; directional
convexity is not assumed. The common framework of the author's paper \emph{Shortest-side
skeletons of edge-perimeter-minimizing polyhexes} is reproduced for self-containment.
\end{abstract}

\section{The hexagonal lattice, runs, and directional convexity}\label{sec:prelim}

This part reproduces, for self-containment, the common framework of the main paper
[Utsunomiya, \emph{Shortest-side skeletons of edge-perimeter-minimizing polyhexes}]:
the lattice, the perimeter functional, directional convexity, the tight hull and its
identities, and the extremal capacity function. Numbering follows this note.

\subsection{Lattice and polyhexes}\label{sec:2.1}

Cells of the hexagonal tiling are indexed by axial coordinates
\((q,r)\in\Z^{2}\). We call \(q\), \(r\) and
\[
   s:=q+r
\]
the three \emph{keys} of a cell; in words, the third key of a cell is the sum of its two
axial coordinates. Two cells are \emph{adjacent} when they differ by one of the six
vectors
\[
   d_{0}=(1,0),\quad d_{1}=(0,1),\quad d_{2}=(-1,1),\quad
   d_{3}=(-1,0),\quad d_{4}=(0,-1),\quad d_{5}=(1,-1).
\]

A \emph{polyhex} is a finite nonempty set of cells that is connected with respect to
adjacency. We write \(|S|\) for the number of cells of \(S\).

The following four facts are immediate from the definitions and are used throughout
without further comment.

\begin{fact}\label{f:1}
Two distinct cells are adjacent if and only if they agree in exactly one key and differ
by exactly \(1\) in the other two. In words: an adjacency step keeps one of the three
keys fixed and changes each of the other two by one unit (the step \(d_{0}=(1,0)\), for
instance, sends \((q,r,s)\) to \((q+1,r,s+1)\), so the two changing keys need not move
in opposite senses).
\end{fact}

\begin{proof}
Each \(d_{j}\) changes \((q,r,s)\) by \((1,0,1)\), \((0,1,1)\), \((-1,1,0)\),
\((-1,0,-1)\), \((0,-1,-1)\) or \((1,-1,0)\), and these are exactly the six triples with
one zero entry and two entries equal to \(\pm1\) that are realized. Conversely a
difference of that shape is one of the \(d_{j}\).
\end{proof}

\begin{fact}\label{f:2}
A \emph{line} is a set of cells on which one key is constant. Two lines belonging to
different keys meet in exactly one cell; two distinct lines belonging to the same key are
disjoint.
\end{fact}

\begin{fact}\label{f:3}
One adjacency step changes each of \(q\), \(r\) and \(s\) by at most \(1\).
\end{fact}

\begin{fact}\label{f:4}
For a polyhex \(S\) and each key, the set of values taken by that key on \(S\) is an
interval of consecutive integers. In words: no key value inside the range of \(S\) can be
skipped.
\end{fact}

\begin{proof}
Let \(k\) be a key and let \(x,y\in S\) realize two values of \(k\). A path of adjacent
cells inside \(S\) joins them, and by Fact~\ref{f:3} the value of \(k\) changes by at most
one per step, so every intermediate integer occurs on the path.
\end{proof}

\subsection{Edge perimeter, runs and directional convexity}\label{sec:2.2}

Let \(a(S)\) denote the number of unordered pairs of adjacent cells of \(S\). The
\emph{edge perimeter} of \(S\) is
\begin{equation}\label{eq:p}
   p(S)=6|S|-2a(S),
\end{equation}
in words: six times the number of cells, minus twice the number of adjacent pairs.
Equivalently, \(p(S)\) counts the cell edges of \(S\) that are not shared with another
cell of \(S\); we call these the \emph{boundary edges}. The \emph{edge semiperimeter} is
\[
   \sigma(K):=\tfrac{1}{2}p(K)
\]
for any finite set of cells \(K\); in words, half the edge perimeter. That this is an
integer is part of Lemma~\ref{lem:runs} below.

A line \(\Lambda\) is \emph{occupied} by \(S\) when \(S\cap\Lambda\neq\emptyset\); we call
\(S\cap\Lambda\) the \emph{trace} of \(S\) on \(\Lambda\). Cells of a line are linearly
ordered (for a line of constant \(q\) by the value of \(r\), for one of constant \(r\) by
the value of \(q\), for one of constant \(s\) by the value of \(q\)), and consecutive
cells in that order are adjacent. A \emph{run} of \(S\) on \(\Lambda\) is a maximal block
of consecutive cells of the trace. Let
\[
   R(S)=\text{the total number of runs of \(S\) over all lines},\qquad
   \lambda(S)=\text{the number of occupied lines}.
\]

\begin{definition}\label{def:dconv}
A polyhex \(S\) is \emph{directionally convex} when the trace of \(S\) on every occupied
line is a single run. In words: along each of the three directions, \(S\) never leaves a
line and comes back to it.
\end{definition}

This is a lattice notion and not Euclidean convexity; it is the notion called
\emph{layer convexity} in parts of the literature on lattice animals.

\begin{lemma}\label{lem:runs}
For every finite set of cells \(S\),
\begin{equation}\label{eq:pruns}
   p(S)=2R(S)\ \ge\ 2\lambda(S),
\end{equation}
in words: the edge perimeter equals twice the total number of runs, and the total number
of runs is at least the number of occupied lines. In particular \(p(S)\) is even and
\(\sigma(S)=R(S)\). For a polyhex, equality \(R(S)=\lambda(S)\) holds if and only if
\(S\) is directionally convex.
\end{lemma}

\begin{proof}
Let \(e=\{x,y\}\) be a boundary edge, with \(x\in S\) and \(y\notin S\) the two cells it
separates. By Fact~\ref{f:1} the cells \(x\) and \(y\) agree in exactly one key, hence
they are consecutive cells of exactly one line \(\Lambda\), and \(x\) is an end of the run
of \(S\) on \(\Lambda\) containing it. Conversely a run has exactly two ends, and the cell
immediately beyond each end lies outside \(S\) and determines a boundary edge; for a run
consisting of a single cell the two ends are the two opposite sides of that cell, so they
give two distinct boundary edges. This is a bijection between boundary edges and ends of
runs, whence \(p(S)=2R(S)\).

Every occupied line carries at least one run, so \(R(S)\ge\lambda(S)\), with equality
exactly when every occupied line carries exactly one run, which is
Definition~\ref{def:dconv}.
\end{proof}

\begin{remark}[edge perimeter versus site perimeter]\label{rem:site}
The \emph{site perimeter} of \(S\) is the number of cells outside \(S\) adjacent to at
least one cell of \(S\). It is a different functional from \eqref{eq:p}, and the two have
different minimizers in general; results about one do not transfer to the other. Nothing
in this paper depends on the site perimeter. We mention it only because part of the
literature we cite in Section~\ref{sec:2.5} is written in those terms.
\end{remark}

\subsection{Extreme support lines and traces}\label{sec:2.3}

Let \(S\) be a polyhex. For each key, Fact~\ref{f:4} makes the occupied values an interval,
so there is a smallest and a largest occupied line in each of the three families. These six
lines are the \emph{extreme lines} of \(S\); we list them in the cyclic order
\[
   \Lambda_{1}:\ q=q_{\min},\quad
   \Lambda_{2}:\ r=r_{\max},\quad
   \Lambda_{3}:\ s=s_{\max},\quad
   \Lambda_{4}:\ q=q_{\max},\quad
   \Lambda_{5}:\ r=r_{\min},\quad
   \Lambda_{6}:\ s=s_{\min},
\]
so that \(\Lambda_{i}\) and \(\Lambda_{i+3}\) are the two extreme lines of one and the same
family, indices being read in \(\Z/6\). If a family occupies a single line then the two
corresponding extreme lines coincide; all statements below are read with that degeneracy
allowed.

\begin{definition}\label{def:traces}
Let \(S\) be a \emph{directionally convex} polyhex. Its six \emph{support traces} are the
traces \(S\cap\Lambda_{i}\), each a single run by Definition~\ref{def:dconv}, and we write
\[
   r_{i}:=|S\cap\Lambda_{i}|\ \ge 1 ,\qquad i\in\Z/6,
   \qquad\text{and}\qquad
   \rho(S):=\min_{i} r_{i} .
\]
In words: \(r_{i}\) is the number of cells of \(S\) on the \(i\)-th extreme line, and
\(\rho(S)\) is the smallest of those six numbers.
\end{definition}

\begin{remark}\label{rem:rho}
Definition~\ref{def:traces} is stated for directionally convex \(S\) only. On that domain \(\rho(S)\) is at once the length of the shortest support trace and the
length of the shortest extreme-line run. For
a general polyhex a trace need not be a run, and we do not extend the definition.
\end{remark}

\section{Tight hulls and exact identities}\label{sec:hulls}

\subsection{Tight lattice hulls}\label{sec:2.4}

\begin{definition}\label{def:hexagon}
Let \(L\ge1\) and \(u,v,w\ge0\) be integers satisfying the \emph{admissibility} conditions
\begin{equation}\label{eq:adm}
   u+v\le L-1,\qquad v+w\le L-1,\qquad u+w\le L-1 ,
\end{equation}
and put \(A:=L-u\), \(B:=L-v\), \(C:=L-w\), so that \(A,B,C\ge1\). The corresponding
\emph{lattice hexagon} is
\[
   H(L;u,v,w):=\{\,(q,r):\ 0\le q\le A-1,\ \ 0\le r\le B-1,\ \ w\le q+r\le L-1\,\}.
\]
In words: the triangle of side \(L\) with three corner triangles of sides \(u\), \(v\) and
\(w\) cut off. Degenerate shapes, in which one or more of the cut corners is empty or one
of the three widths equals \(1\), are allowed.
\end{definition}

\begin{definition}\label{def:hull}
The \emph{tight hull} \(H(S)\) of a polyhex \(S\) is the lattice hexagon determined by the
extreme values of the three keys on \(S\). Concretely, normalize by a translation so that
\(q_{\min}=r_{\min}=0\) and set
\[
   A=1+q_{\max},\quad B=1+r_{\max},\quad L=1+s_{\max},\quad w=s_{\min},\quad
   u=L-A,\quad v=L-B .
\]
\end{definition}

\begin{lemma}[the tight hull is admissible]\label{lem:tightadm}
For every polyhex \(S\), the parameters \(L,u,v,w\) of Definition~\ref{def:hull} satisfy
\(L\ge1\), \(u,v,w\ge0\) and the three inequalities \eqref{eq:adm}, so
\(H(S)=H(L;u,v,w)\) is a lattice hexagon.
\end{lemma}

\begin{proof}
After the normalization every cell has \(q,r\ge0\), hence \(s\ge0\), so \(w\ge0\) and
\(L\ge1\); and \(u=s_{\max}-q_{\max}\), \(v=s_{\max}-r_{\max}\). A cell with
\(q=q_{\max}\) has \(s\ge q_{\max}\), so \(u\ge0\); likewise \(v\ge0\). A cell with
\(s=s_{\max}\) has \(s_{\max}=q+r\le q_{\max}+r_{\max}\), which is \(u+v\le L-1\). A cell
with \(q=0\) has \(s=r\le r_{\max}\), so \(s_{\min}\le r_{\max}\), which is \(v+w\le L-1\);
a cell with \(r=0\) gives \(s_{\min}\le q_{\max}\), which is \(u+w\le L-1\).
\end{proof}

For a lattice hexagon \(H\) we write \(\Lambda_{1},\dots,\Lambda_{6}\) for its six extreme
lines in the cyclic order of Section~\ref{sec:2.3}, namely
\[
   q=0,\quad r=B-1,\quad s=L-1,\quad q=A-1,\quad r=0,\quad s=w ,
\]
and call them the \emph{side lines}; we write
\[
   \ell_{i}:=|H\cap\Lambda_{i}|
\]
for the six \emph{side lengths}, \(h:=\min_{i}\ell_{i}\) for the shortest side length, and
\(\kappa_{i}:=H\cap\Lambda_{i}\cap\Lambda_{i+1}\) for the \emph{corner} cells. Finally we
set
\[
   \delta(H):=\ell_{1}-\ell_{4},
   \qquad
   Q(H):=(\ell_{1}-\ell_{3})^{2}+(\ell_{3}-\ell_{5})^{2}+(\ell_{5}-\ell_{1})^{2} .
\]

\begin{lemma}\label{lem:hull}
Let \(S\) be a polyhex and \(H=H(S)\). Then
\begin{enumerate}
\item[(i)] \(S\subseteq H\);
\item[(ii)] \(S\) and \(H\) occupy exactly the same lines, so \(\lambda(S)=\lambda(H)\);
\item[(iii)] each of the six side lines of \(H\) meets \(S\).
\end{enumerate}
\end{lemma}

\begin{proof}
(i) is immediate from Definition~\ref{def:hull}, since every cell of \(S\) satisfies the
six inequalities defining \(H\). For (ii), the lines occupied by \(S\) in the \(q\)-family
are, by Fact~\ref{f:4}, exactly \(q=0,\dots,A-1\), and likewise \(r=0,\dots,B-1\) and
\(s=w,\dots,L-1\); these are precisely the lines occupied by \(H\), as read off from
Definition~\ref{def:hexagon}. Assertion (iii) is the case of (ii) for the six extreme
values: each of them is attained by some cell of \(S\).
\end{proof}

Lemma~\ref{lem:hull} is the only place where connectedness of \(S\) is used in this
section, and it is used through Fact~\ref{f:4}: for a disconnected set the occupied key
values need not form an interval, and (ii) fails.

\subsection{Hull identities}\label{sec:2.5}

\begin{lemma}\label{lem:hex}
Let \(H=H(L;u,v,w)\) be a lattice hexagon. Then:
\begin{enumerate}
\item[(i)] the six side lengths are
\[
   \ell_{1}=L-v-w,\quad \ell_{2}=v+1,\quad \ell_{3}=L-u-v,\quad
   \ell_{4}=u+1,\quad \ell_{5}=L-u-w,\quad \ell_{6}=w+1 ,
\]
and each side line meets \(H\) in a single run whose two end cells are the corners
\(\kappa_{i-1}\) and \(\kappa_{i}\);
\item[(ii)] \(H\) is a directionally convex polyhex with
\[
   \lambda(H)=\sigma(H)=A+B+C=3L-(u+v+w);
\]
\item[(iii)] the six side lengths satisfy
\begin{equation}\label{eq:sidesum}
   \sum_{i=1}^{6}\ell_{i}=\lambda(H)+3 ,
\end{equation}
in words: the six side lengths of the hull add up to three more than the number of lines
the hull occupies; and the three differences across opposite sides agree,
\begin{equation}\label{eq:delta}
   \ell_{1}-\ell_{4}=\ell_{3}-\ell_{6}=\ell_{5}-\ell_{2}=\delta(H)=L-u-v-w-1 ;
\end{equation}
\item[(iv)] with \(T(k):=k(k+1)/2\),
\[
   |H|=T(L)-T(u)-T(v)-T(w) ;
\]
\item[(v)] \textup{(hull defect identity)}
\begin{equation}\label{eq:defectid}
   \lambda(H)^{2}+3-12|H| \;=\; 3\,\delta(H)^{2}+2\,Q(H)\;\ge\;0 ,
\end{equation}
in words: the square of the number of occupied lines, plus three, minus twelve times the
number of cells, equals three times the square of the common opposite-side difference plus
twice the imbalance \(Q(H)\).
\end{enumerate}
\end{lemma}

\begin{proof}
(i) On the line \(q=0\) the defining inequalities reduce to \(w\le r\le\min(B-1,L-1)=B-1\),
which is a nonempty interval because \(w\le B-1\) is the condition \(v+w\le L-1\) of
\eqref{eq:adm}; its length is \(B-w=L-v-w\). On the line \(r=B-1\) they reduce to
\(\max(0,w-B+1)\le q\le\min(A-1,v)\), and \eqref{eq:adm} gives \(w-B+1\le0\) and
\(v\le A-1\), so the interval is \(0\le q\le v\), of length \(v+1\). The remaining four
sides are obtained in the same way, each time using exactly one of the three inequalities
of \eqref{eq:adm} at each end. Each side is an interval of consecutive cells of its line,
hence a run; its two end cells are the two cells where a second defining inequality becomes
an equality, that is, the two corners.

(ii) The occupied lines are \(q=0,\dots,A-1\), \(r=0,\dots,B-1\) and \(s=w,\dots,L-1\),
which number \(A+B+C\); on each of them the trace is an interval of consecutive cells by
the same computation as in (i), so \(H\) is directionally convex and
Lemma~\ref{lem:runs} gives \(\sigma(H)=R(H)=\lambda(H)\). Connectedness follows since each
trace is a run and consecutive traces meet a common line.

(iii) Summing the six expressions of (i),
\[
   \sum_{i}\ell_{i}=3L-(u+v+w)+3=\lambda(H)+3 .
\]
Using the side-length formulas of (i), the three opposite-side differences
\(\ell_{1}-\ell_{4}\), \(\ell_{3}-\ell_{6}\), \(\ell_{5}-\ell_{2}\) of \eqref{eq:delta}
are \((L-v-w)-(u+1)\), \((L-u-v)-(w+1)\) and \((L-u-w)-(v+1)\), all equal to
\(L-u-v-w-1\).

(iv) The triangle \(\{q,r\ge0,\ q+r\le L-1\}\) has \(T(L)\) cells. The three sets removed
in Definition~\ref{def:hexagon} are the corner triangles \(q\ge A\), \(r\ge B\) and
\(q+r\le w-1\), of \(T(u)\), \(T(v)\) and \(T(w)\) cells respectively, and \eqref{eq:adm}
makes them pairwise disjoint: a cell with \(q\ge A\) and \(r\ge B\) would have
\(q+r\ge A+B=2L-u-v\ge L+1\), which is impossible, and a cell with \(q\ge A\) and
\(q+r\le w-1\) would give \(r\le w-1-A<0\), and symmetrically.

(v) Put \(U:=u+v+w\) and \(P:=u^{2}+v^{2}+w^{2}\). By (ii) and (iv),
\[
   \lambda(H)^{2}+3-12|H|=(3L-U)^{2}+3-6L(L+1)+6P+6U
   =3L^{2}-6LU-6L+U^{2}+6P+6U+3 .
\]
On the other side, \(\delta(H)=L-U-1\) and, by (i),
\(Q(H)=(u-w)^{2}+(w-v)^{2}+(v-u)^{2}=3P-U^{2}\), so
\[
   3\delta(H)^{2}+2Q(H)=3(L-U-1)^{2}+2(3P-U^{2})
\]
expands to the same expression. In words: both sides equal three times \(L^{2}\), minus six
times \(LU\), minus six times \(L\), plus \(U^{2}\), plus six times \(P\), plus six times
\(U\), plus three. Nonnegativity is clear from the right-hand side.
\end{proof}



\begin{remark}[attribution]\label{rem:fs}
An equivalent defect identity is obtained from F\"ulep and Sieben's polyiamond identity
\cite[Lemma~3.3]{FulepSieben2010} through the triangular--hexagonal duality developed in their Section~5.
In the same site-perimeter variables, \cite[Proposition~5.10]{FulepSieben2010} yields exactly the capacity
function \(M(t)\) of Section~\ref{sec:2.6} after the substitution \(p=t+3\); the value
\(p=7\) excluded there corresponds to the value \(t=4\) excluded in
Remark~\ref{rem:attain} below. We nevertheless give the derivation here in edge-perimeter
variables to keep the argument self-contained, and we do not identify the two settings:
\cite{FulepSieben2010} is about the site perimeter (Remark~\ref{rem:site}).
\end{remark}

\section{Capacity and minimum edge perimeter}\label{sec:capacity}

\subsection{Extremal capacity}\label{sec:2.6}

We first record the two operations on directionally convex polyhexes that the capacity
computation needs: removing a whole support trace, and putting it back one cell at a time.

\begin{lemma}[peeling]\label{lem:peel}
Let \(K\) be a directionally convex polyhex, let \(\Lambda\) be one of its six extreme
lines, and let \(K':=K\sm\Lambda\). If \(K'\neq\emptyset\), then
\begin{enumerate}
\item[(i)] \(K'\) is a directionally convex polyhex with \(\sigma(K')\le\sigma(K)-1\);
\item[(ii)] the cells of \(K\cap\Lambda\) can be ordered \(c_{1},\dots,c_{k}\) so that for
every \(j\) the set \(K'\cup\{c_{1},\dots,c_{j}\}\) is a directionally convex polyhex of
edge semiperimeter at most \(\sigma(K)\).
\end{enumerate}
\end{lemma}

\begin{proof}
Say \(\Lambda\) is the line \(q=q_{\min}\); the other five cases are the same after
permuting the roles of the keys and of the two senses.

(i) For each occupied value \(k\) of \(q\) let \(K_{k}\) be the trace of \(K\) on
\(q=k\); it is a run. If \(K\) is connected, then for consecutive values \(k,k+1\) a path
inside \(K\) joining a cell of \(K_{k}\) to a cell of \(K_{k+1}\) must, by
Fact~\ref{f:3}, contain a step from a cell with \(q=k\) to a cell with \(q=k+1\), so
\(K_{k}\) and \(K_{k+1}\) are adjacent. Hence \(K'\), the union of the \(K_{k}\) with
\(k\ge q_{\min}+1\), is connected. Its traces in the \(q\)-family are traces of \(K\),
hence runs. Now fix a line of the \(r\)-family, \(r=k\); its cells are \((q,k)\), so the
trace of \(K\) on it is a run if and only if its \(q\)-values form an interval, and by
Fact~\ref{f:2} it meets \(\Lambda\) in at most the one cell \((q_{\min},k)\). The trace of
\(K'\) on this line is the trace of \(K\) with that cell, if present, removed. Since
\(q_{\min}\) is the least \(q\)-value occurring anywhere in \(K\), it is the least
\(q\)-value of the interval, so removing it leaves an interval, possibly empty. A line of
the \(s\)-family, \(s=k\), has cells \((q,k-q)\), again parametrized by \(q\), and the
same sentence applies verbatim. So all traces of \(K'\) are runs and \(K'\) is
directionally convex. Finally \(\Lambda\) is occupied by
\(K\) and not by \(K'\), so \(\lambda(K')\le\lambda(K)-1\), and by Lemma~\ref{lem:runs}
\(\sigma(K')=\lambda(K')\le\lambda(K)-1=\sigma(K)-1\).

(ii) Since \(K\) is connected there is a cell \(c_{1}\in K\cap\Lambda\) adjacent to
\(K'\). Order the remaining cells of the run \(K\cap\Lambda\) so that every initial
segment \(\{c_{1},\dots,c_{j}\}\) is a block of consecutive cells of that run containing
\(c_{1}\). Write \(K_{j}:=K'\cup\{c_{1},\dots,c_{j}\}\). Then \(K_{j}\) is connected,
because \(c_{1}\) is adjacent to \(K'\) and each later \(c_{j}\) is adjacent to a cell
already inserted. Its trace on \(\Lambda\) is a block, hence a run; its trace on any other
line of the \(q\)-family is a trace of \(K'\), hence a run. On a line of the \(r\)- or
\(s\)-family, parametrized by \(q\) as in (i), the trace of \(K_{j}\) is the trace of
\(K'\) there together with at most one inserted cell, which lies on \(\Lambda\) and so
has \(q=q_{\min}\). The trace of \(K\) on that line is a run containing this cell and, by
(i), the trace of \(K'\) is that run with its least element \(q_{\min}\) removed; putting
the cell back restores the run of \(K\). Hence every trace of \(K_{j}\) is a run, so
\(K_{j}\) is directionally convex, and every line it occupies is occupied by \(K\), whence
\(\sigma(K_{j})=\lambda(K_{j})\le\lambda(K)=\sigma(K)\).
\end{proof}

\begin{corollary}\label{cor:subsets}
Let \(K\) be a directionally convex polyhex. For every \(j\) with \(1\le j\le|K|\) there is
a directionally convex polyhex \(K_{j}\subseteq K\) with \(|K_{j}|=j\) and
\(\sigma(K_{j})\le\sigma(K)\).
\end{corollary}

\begin{proof}
Strong induction on \(|K|\). For \(|K|=1\) there is nothing to prove. Let \(|K|\ge2\). Some
extreme line \(\Lambda\) does not contain all of \(K\): otherwise \(K\) would be contained
in two extreme lines belonging to different families, which by Fact~\ref{f:2} meet in a
single cell, contradicting \(|K|\ge2\). Apply Lemma~\ref{lem:peel} to that \(\Lambda\). The
sets \(K_{j}\) of part (ii) supply all sizes between \(|K'|\) and \(|K|\); the sizes below
\(|K'|\) are supplied by the inductive hypothesis applied to \(K'\), whose edge
semiperimeter is at most \(\sigma(K)\).
\end{proof}

\begin{definition}\label{def:M}
For an integer \(t\ge0\) put
\[
   M(t):=\left\lfloor\frac{t^{2}+3}{12}\right\rfloor ,
\]
in words: the integer part of one twelfth of the square of \(t\) plus three. Thus
\(M(0)=M(1)=M(2)=0\) and \(M(3)=1\). The function \(M\) is nondecreasing.
\end{definition}

\begin{theorem}\label{thm:capacity}
Let \(N\ge1\).
\begin{enumerate}
\item[(a)] Every polyhex \(S\) satisfies \(|S|\le M(\sigma(S))\); in words, a configuration
of edge semiperimeter \(t\) cannot have more than \(M(t)\) cells.
\item[(b)] For every \(t\ge3\) there is a directionally convex polyhex \(K\) with
\(|K|=M(t)\) and \(\sigma(K)\le t\).
\item[(c)] The minimum edge perimeter of an \(N\)-cell polyhex is
\[
   p_{\min}(N)=2g(N),\qquad g(N)=\left\lceil\sqrt{12N-3}\,\right\rceil ,
\]
in words: the smallest edge semiperimeter that can hold \(N\) cells is the least integer at
least as large as the square root of twelve times \(N\) minus three, and the minimum edge
perimeter is twice that number.
\item[(d)] Every \(N\)-cell minimizer \(S\) is directionally convex and satisfies
\(\sigma(S)=\lambda(S)=g(N)\).
\item[(e)] Writing \(\exc(S):=p(S)-p_{\min}(|S|)\), a polyhex \(S\) satisfies
\(\exc(S)\le2\) if and only if \(|S|>M(\sigma(S)-2)\).
\end{enumerate}
\end{theorem}

The directional convexity in (d) is a known result: Davoli, Piovano and Stefanelli
\cite[Proposition~3.3]{DavoliPiovanoStefanelli2017} prove that every minimizer of the
edge-isoperimetric problem on the triangular lattice is \(3\)-convex, which is
directional convexity in the present terminology; the proof below is the self-contained
one of the main paper, via the hull identities.

\begin{proof}
(a) Let \(H=H(S)\). By Lemma~\ref{lem:hull} and Lemma~\ref{lem:hex}(v),
\[
   12|S|\le12|H|=\lambda(H)^{2}+3-\bigl(3\delta(H)^{2}+2Q(H)\bigr)\le\lambda(S)^{2}+3
   \le\sigma(S)^{2}+3 ,
\]
the last step by Lemma~\ref{lem:runs}. Since \(|S|\) is an integer, \(|S|\le M(\sigma(S))\).

(b) For \(t\equiv3\pmod 6\), say \(t=6m+3\), take \(K=H(3m+1;m,m,m)\); for \(t=6m\) take
\(H(3m-1;m-1,m-1,m-1)\); for \(t=6m+1\) take \(H(3m;m,m,m-1)\); for \(t=6m+2\) take
\(H(3m+1;m,m,m+1)\); for \(t=6m+4\) take \(H(3m+1;m,m,m-1)\); for \(t=6m+5\) take
\(H(3m+2;m,m,m+1)\). In each case \eqref{eq:adm} holds for the indicated range of \(m\),
Lemma~\ref{lem:hex}(ii) gives \(\sigma(K)=t\), and Lemma~\ref{lem:hex}(iv),(v) give
\(12|K|=t^{2}+3-D\) with \(D\in\{0,3,4,7\}\) according to the residue of \(t\), so
\(|K|=M(t)\). The single case not covered by an admissible triple is \(t=4\); there
\(M(4)=1\) and a single cell has edge semiperimeter \(3\le4\).

(c) Put \(t_{N}:=\min\{t\ge3:\ M(t)\ge N\}\). By (a) every \(N\)-cell polyhex has
\(M(\sigma(S))\ge N\), hence \(\sigma(S)\ge t_{N}\). Conversely (b) produces a
directionally convex \(K\) with \(|K|=M(t_{N})\ge N\) and \(\sigma(K)\le t_{N}\), and
Corollary~\ref{cor:subsets} extracts from it an \(N\)-cell polyhex of edge semiperimeter at
most \(t_{N}\). Therefore \(p_{\min}(N)=2t_{N}\). Finally, for integers \(t\),
\[
   M(t)\ge N\iff \frac{t^{2}+3}{12}\ge N\iff t\ge\sqrt{12N-3},
\]
and \(\sqrt{12N-3}\ge3\) for \(N\ge1\), so \(t_{N}=g(N)\).

(d) By (c), \(\sigma(S)=g(N)\). Let \(H=H(S)\). Applying (a) to \(H\), which is
directionally convex with \(\sigma(H)=\lambda(H)\), gives \(M(\lambda(H))\ge|H|\ge N\) and
therefore \(\lambda(H)\ge t_{N}=g(N)\). Combining with Lemmas~\ref{lem:runs}
and~\ref{lem:hull},
\[
   g(N)=\sigma(S)=R(S)\ \ge\ \lambda(S)=\lambda(H)\ \ge\ g(N),
\]
so \(R(S)=\lambda(S)\) and \(S\) is directionally convex.

(e) Write \(t=\sigma(S)\) and \(k=|S|\); by (c), \(\exc(S)=2\bigl(t-g(k)\bigr)\). Since
both sides are integers, \(\exc(S)>2\) is equivalent to \(g(k)\le t-2\). If \(t-2\ge3\),
then by the definition of \(g\) as \(t_{k}\) and the monotonicity of \(M\) this is
equivalent to \(M(t-2)\ge k\). If \(t-2<3\), then \(t\in\{3,4\}\) because a polyhex has
\(t\ge3\); here \(M(t-2)=0<k\) and \(g(k)\ge3>t-2\), so both statements are false.
In both cases \(\exc(S)>2\) is equivalent to \(M(t-2)\ge k\), which is the assertion.
\end{proof}

\begin{remark}\label{rem:attain}
Parts (a) and (b) together say exactly that
\[
   M(t)=\max\{\,|K|:\ K\text{ a polyhex with }\sigma(K)\le t\,\},
\]
in words: \(M(t)\) is the largest number of cells that fit within an edge-semiperimeter
budget of \(t\). The budget form is the one the rest of the paper uses, and the
inequality in (b) is not a weakening: for every \(t\ge3\) other than \(t=4\) the
hexagons listed in the proof do satisfy \(\sigma(K)=t\), and the value \(t=4\) is not the
edge semiperimeter of any polyhex at all. Indeed a polyhex \(K\) with \(\sigma(K)=4\)
would have \(|K|\le M(4)=1\) by (a), while the one-cell polyhex has edge semiperimeter
\(3\).
\end{remark}

\subsection{The plateau bridge}\label{sec:2.7}

The function \(N\mapsto g(N)\) is nondecreasing and takes each of its values on a block of
consecutive integers. For \(\gamma\) in the range of \(g\) let
\[
   N_{b}(\gamma):=\min\{\,N\ge1:\ g(N)=\gamma\,\}
\]
be the first cell count at which the minimum edge semiperimeter equals \(\gamma\).

\begin{lemma}\label{lem:bridge}
For every \(\gamma\) in the range of \(g\),
\begin{equation}\label{eq:bridge}
   N_{b}(\gamma)=M(\gamma-1)+1 ,
\end{equation}
in words: the plateau on which the minimum edge semiperimeter equals \(\gamma\) begins one
cell after the capacity of edge semiperimeter \(\gamma-1\) has been exhausted.
\end{lemma}

\begin{proof}
For integers \(N\ge1\) and \(t\ge3\) we have, as in the proof of
Theorem~\ref{thm:capacity}(c), \(g(N)\le t\) if and only if \(N\le M(t)\). Hence
\(\{N:\ g(N)\le\gamma-1\}=\{1,\dots,M(\gamma-1)\}\), so the least \(N\) with
\(g(N)\ge\gamma\) is \(M(\gamma-1)+1\); as \(\gamma\) lies in the range of \(g\), that
least \(N\) satisfies \(g(N)=\gamma\).
\end{proof}




\section{Two lemmas}

\subsection{A degree-two articulation cell forces excess at least four}\label{sec:11.3}


\begin{lemma}[arithmetic lemma]\label{lem:arith}
Let \(a,b\ge1\) be integers with \((a,b)\neq(1,1)\). Then
\(g(a)+g(b)\ge g(a+b+1)+1\). For \((a,b)=(1,1)\) the statement fails.
\end{lemma}

\begin{proof}
Put \(A:=g(a)\), \(B:=g(b)\). By definition \(A\ge\sqrt{12a-3}\), that is
\(12a\le A^{2}+3\), and likewise for \(b\); hence
\(12(a+b+1)-3\le A^{2}+B^{2}+15\), and by monotonicity of the ceiling of the square root,
\(g(a+b+1)\le\lceil\sqrt{A^{2}+B^{2}+15}\,\rceil\). Since \(A+B-1\) is a positive integer,
\(\lceil\sqrt{A^{2}+B^{2}+15}\,\rceil\le A+B-1\) is equivalent to
\(A^{2}+B^{2}+15\le(A+B-1)^{2}\), that is to \((A-1)(B-1)\ge8\). Now \(g\) is nondecreasing
with \(g(1)=3\) and \(g(2)=5\), so \(A=3\) if \(a=1\) and \(A\ge5\) if \(a\ge2\), and
similarly for \(B\); unless \(a=b=1\) we get \((A-1)(B-1)\ge2\cdot4=8\). Finally
\(A=B=3\) happens exactly for \(a=b=1\), where \(g(1)+g(1)=6=g(3)\).
\end{proof}

\begin{lemma}[Bridge-Excess Lemma]\label{lem:bridge-exc}
Let \(X\) be a polyhex with \(|X|=N\ge4\) having a degree-two articulation cell. Then
\(\exc(X)\ge4\).
\end{lemma}

\begin{proof}
Let \(c\) be such a cell. Every component of \(X\sm\{c\}\) contains a neighbour of \(c\),
and \(c\) has two neighbours, so there are at most two components; being an articulation
cell there are at least two. Call them \(X_{1},X_{2}\) and put \(n_{i}=|X_{i}|\), so
\(n_{1}+n_{2}=N-1\). Counting adjacent pairs, \(a(X)=a(X_{1})+a(X_{2})+2\), whence
\[
   p(X)=6N-2a(X)=p(X_{1})+p(X_{2})+2 .
\]
In words: the two pieces keep their own perimeters and the bridge contributes two further
boundary edges. With \(p(X_{i})\ge2g(n_{i})\) and Lemma~\ref{lem:arith},
\[
   p(X)\ \ge\ 2g(n_{1})+2g(n_{2})+2\ \ge\ 2\bigl(g(N)+1\bigr)+2\ =\ p_{\min}(N)+4 .
\]
Lemma~\ref{lem:arith} needs \((n_{1},n_{2})\neq(1,1)\), which would mean \(N=3\).
\end{proof}


\section{Extreme endpoints are never articulation cells}\label{sec:15.4}

In the one-cell deletion proof of the growth theorem (recorded in a companion note) one
needs that \emph{some} endpoint of a shortest extreme line of a minimizer can be deleted
without disconnecting the configuration. The local argument used there does not in fact
depend on the line being shortest, nor on directional convexity, nor on which end is taken;
it depends only on the excess bound. Making this explicit gives the following statement.

For a general polyhex \(S\), the six extreme lines are defined as in
Section~\ref{sec:2.3}, and an \emph{endpoint} of the trace of \(S\) on one of them is a
cell of the trace with extremal position along that line. When \(S\) is directionally
convex these are the two ends of the support trace.

\begin{theorem}\label{thm:endpoint}
Let \(S\) be a polyhex with \(|S|\ge4\) and \(\exc(S)\le2\). Then no endpoint of any extreme
line of \(S\) is an articulation cell of \(S\).
\end{theorem}

\begin{proof}
Let \(c\) be an endpoint of the trace of \(S\) on an extreme line \(\Lambda\). The symmetry
group of the lattice contains the rotations through \(60^{\circ}\), which permute the six
extreme roles cyclically, and the reflection \((q,r)\mapsto(q,-q-r)\), which fixes every
line of the \(q\)-family and reverses the order of the cells along it; both preserve
adjacency, hence \(p\), \(|S|\) and \(\exc\). Applying a suitable composition of these
symmetries, we may assume that \(\Lambda\) is the line \(q=q_{\max}\) and that \(c\) has
the largest value of \(r\) among the cells of \(S\) on it.

The directions \(d_{0}=(1,0)\) and \(d_{5}=(1,-1)\) increase \(q\), and \(c\) lies on
\(q=q_{\max}\); the direction \(d_{1}=(0,1)\) preserves \(q\) and increases \(r\), and
\(c\) has the largest \(r\) on that line. So \(c+d_{0}\), \(c+d_{1}\), \(c+d_{5}\) are not
in \(S\), and the neighbours of \(c\) in \(S\) lie in the three consecutive directions
\(d_{2},d_{3},d_{4}\). If \(\dg_{S}(c)\le1\) then \(S\sm\{c\}\) is connected. If
\(\dg_{S}(c)=3\), or if \(\dg_{S}(c)=2\) with neighbours in two adjacent directions, then
those neighbours are themselves adjacent (for instance \(c+d_{2}\) and \(c+d_{3}\) differ
by \(d_{1}\)), so any path through \(c\) can be rerouted and \(c\) is not an articulation
cell. Hence if \(c\) were an articulation cell then \(\dg_{S}(c)=2\) with neighbours in the
directions \(d_{2}\) and \(d_{4}\). Since \(|S|\ge4\), the Bridge-Excess
Lemma~\ref{lem:bridge-exc} would then give \(\exc(S)\ge4\), contradicting
\(\exc(S)\le2\).
\end{proof}

In words: a cell at the end of an extreme line has all its neighbours on one side, in three
consecutive directions, so it can only separate the configuration by being a degree-two
bridge; and a degree-two bridge costs at least four units of edge perimeter.

Both hypotheses are needed. In the bent three-cell polyhex \(\{(0,1),(0,2),(1,0)\}\), of
excess \(2\), the cell \((0,1)\) is an endpoint of the extreme line \(q=0\) and its deletion
disconnects the set; and in the four-cell polyhex \(\{(0,1),(0,2),(0,3),(1,0)\}\), of excess
\(4\), the same is true of \((0,1)\). Directional convexity is not needed: the nine-cell
configuration
\[
   \{(0,1),(0,2),(0,3),(1,0),(1,1),(1,2),(1,3),(2,0),(2,2)\}
\]
has excess \(2\) and is not directionally convex, and the theorem applies to it.

\bibliographystyle{plain}
\bibliography{refs}
\end{document}
